\section{Tasks and Evaluation}
\label{sec:tasks}

% Original scaffold TODOs: task inputs, outputs, capabilities, construction, metrics, ambiguity.
We propose two core tasks and one exploratory relation task.
\todo{Confirm the final task set and freeze the annotation schema before test construction.}

\subsection{Cultural Cue Recognition}
\label{sec:cue-task}
\paragraph{Input and output.}
The input is a video unit with the modalities specified by its evaluation condition.
The output is a set of records $\{(c_i,A_i,k_i)\}$, where $c_i$ is a cue concept, $A_i$ is its set of cultural associations, and $k_i$ is its component type.
Cue recognition pools the supported associations of repeated occurrences
of a concept within an evaluation unit. The relation task uses separate
instance identifiers.
The main condition does not supply the expected cues, a search query, or an uploader title.
The model may return an empty set or state that the available evidence does not support a more specific association.

\paragraph{Ability measured.}
The task tests whether the model preserves distinct cultural associations within one scene or event.
A correct event location does not substitute for a correct association of each practice or artifact.
An association can be shared by several communities and should not be forced into a single-country answer.

\paragraph{Construction and scoring.}
The proposed reference uses canonical cue identifiers with accepted aliases.
Predicted and reference cue records are aligned one-to-one by concept identifier; repeated identical records count as extra predictions.
Cue precision, recall, and F1 measure discovery.
Association-set F1 and component accuracy are reported on matched cues, alongside end-to-end record F1 requiring the concept, association set, and component to be correct.
This prevents conditional scores from concealing missed cues.
We also report macro averages by cultural association and domain with the corresponding sample counts.
\todo{Confirm concept-level aggregation versus instance-level recognition; define aliases, acceptable association sets, and adjudication of open-vocabulary answers.}

\subsection{Cue-Modality Attribution}
\label{sec:modality-task}
\paragraph{Input and output.}
The model receives a video and a validated cue identifier or short description.
It predicts the set of evidence channels that support that cue in the supplied input.
An empty set is valid when the supplied condition contains no supporting evidence.
The candidate channels are visual appearance, visual action, music or other non-speech sound, speech, and visible text.
An externally supplied subtitle or transcript has its own label when that input is available.
\todo{Validate channel boundaries and whether appearance/action should be separate subclasses of the visual modality.}

\paragraph{Ability measured.}
The task assesses attribution of an interpretation to observable evidence.
The presence of music alone does not establish that music supports the cultural label.
Annotators must identify a relevant passage, action, object, or sound and explain its connection to the cue.
A textual assertion inside a video is distinguished from a direct depiction of the practice.
Providing the cue isolates evidence attribution from cue-discovery errors.

\paragraph{Scoring and controls.}
The proposed metrics are micro-F1, macro-F1, and exact-set match over channel labels.
Gold evidence records support auditing of these labels, but precise temporal localization is not required by the core task.
A cue-description-only baseline tests whether the answer can be guessed from the concept rather than the video.
Matched cases in which the same cue is supported by different channels provide a stronger test.
\todo{Specify evidence granularity and the manual procedure for determining that a channel supports a cue.}

\subsection{Candidate: Relation Frames and Roles}
\label{sec:relation-task}
\paragraph{Input and output.}
The initial pairwise setting supplies a video and two validated target nodes in random order.
A node may be a cultural cue or an identified entity, event, or context participating in its realization.
Any cultural association of a carrier or context requires its own evidence.
The model outputs a relation frame and an assignment of node identifiers to that frame's roles.
For symmetric frames, argument order has no effect on correctness.
An insufficient-evidence answer is allowed when the input cannot establish the relation.
This is a classification task over situated elements, not over country names.
\todo{Confirm the node types and whether relations with more than two participants require an explicit result node or an event-level record.}

\paragraph{Candidate frames.}
\begin{description}
    \item[Co-presence.] Two \emph{participants} share a coherent scene or event. This relation alone makes no claim about their integration.
    \item[Embedding.] An \emph{element} occurs within a receiving \emph{context}, such as an event, institution, or place.
    \item[Realization.] A cultural \emph{form} is performed, worn, made, or enacted by a \emph{carrier}.
    \item[Incorporation.] A \emph{component} becomes part of an identifiable \emph{base} artifact or practice.
    \item[Derivation.] A \emph{source} work or repertoire is expressed through a \emph{derivative} realization.
    \item[Synthesis.] Two \emph{contributors} form an integrated result without a stable component--base distinction. The pairwise setting treats the result as event context.
\end{description}
Co-presence and the contributor roles of synthesis are symmetric.
The other frames have different role slots.
For example, incorporation assigns a textile to the component role and a garment to the base role; reversing these assignments changes the claim.
The ordering is a representation convention, not a measure of dominance or historical transfer.
% Conflict from prior discussion: a five-frame option merges realization into embedding.
\todo{Test six frames against the five-frame alternative that merges embedding and realization; confirm whether synthesis can be annotated reliably.}

\paragraph{What counts as a gold relation.}
Multiple relations may hold for one pair or one video.
The annotation should retain all supported frames rather than force one primary label.
Under this working convention, co-presence may accompany a more specific
frame. Established coexistence is distinct from insufficient evidence to
identify an event or its participants.
For an initial diagnostic, we propose a target-frame decision followed by role binding.
Joint prediction can then return a set of supported frame--role records.
Derivation requires evidence of a source relationship; resemblance alone is insufficient.
Cases requiring external provenance belong in a separately specified metadata condition.
\todo{Choose between a set-valued joint task and a curated single-relation subset; decide whether co-presence is independent or only a fallback, and define ambiguity and multiple valid answers.}

\paragraph{Diagnostic settings.}
In \emph{oracle-frame role binding}, the gold frame and target nodes are provided, and the model fills the roles.
In \emph{joint frame and role classification}, target nodes are provided and both the frame and roles are predicted.
The first setting separates role confusion from frame confusion.
Node descriptions must not state the target relation.
Because semantic types such as ``troupe'' and ``dance'' can still reveal roles, a node-text-only baseline is essential.
\todo{Construct balanced node pairs and matched video alternatives to measure gains beyond description-only prediction.}

\paragraph{Metrics.}
For a single target relation, joint exact match requires the frame and every role to be correct.
Frame macro-F1 and oracle-frame role accuracy provide component scores.
If multiple relations are allowed, record precision, recall, and F1 are computed after one-to-one matching, with exact set match reported separately.
Symmetric role permutations are treated as equivalent.
Swap consistency checks that changing node order preserves the semantic prediction.
It is reported alongside accuracy because a consistently wrong model can still pass this check.
\todo{Freeze the primary relation metric after the annotation pilot and specify the treatment of abstention.}

\subsection{Optional Tasks and Evaluation Protocols}
\label{sec:optional-tasks}
Temporal evidence grounding would return a set of intervals for a named cue, extending the localization setting of Cultural Moment Benchmark \citep{satar-etal-2026-cultural}.
Repeated appearances should be separate spans, not one interval covering unrelated content.
One-to-one span matching at a stated temporal IoU threshold supports precision, recall, and F1.
Spatial grounding would be limited to cues with localizable visible regions, following the recognition--grounding distinction in Seeing Culture \citep{satar-etal-2025-seeing}.
\todo{Decide whether to annotate either grounding task and set its evaluation thresholds.}

Cross-language consistency would use validated translations of the same question about the same video, informed by multilingual cultural QA and aligned evaluation \citep{romero2024cvqa,zheng-etal-2026-mmac}.
Agreement and joint correctness should both be reported.
Context robustness would compare matched instances of a practice in different contexts.
Natural recordings can support matched comparisons, but do not isolate the causal effect of context if performers, camera views, and actions also differ.
\todo{Select at most one primary diagnostic extension and establish matched-pair controls.}

Compositional generalization is a split strategy rather than an additional annotation task.
A held-out culture pair or cue--frame combination should use components that each appear in development data.
All versions of one recording and related event must remain in the same split.
For pretrained models, ``unseen'' refers only to the benchmark's declared partitions; it does not establish absence from pretraining.
\todo{Choose a test-only or train/development/test release and verify support for each held-out composition.}
